\section{How to approach a case}
\begin{quotation}
    \textbf{Author's note:} This chapter is based on the first part of the lesson regarding first lab
\end{quotation}

\noindent
The primary goal of a forensic investigation is to produce report that can be evidence in legal proceedings.\\
Unlike technical report, legal report must be understandable to a judge who may not have technical knowledge.
\bigskip

\noindent
A forensic report must possess specific legal and scientific qualities that separate it from techincal documentation:
\begin{itemize}
    \item \textbf{Verifiable:} Every assertion made within the document must be directly traceable to objective, uninterpreted data extracted during the investigation.
    \item \textbf{Replicable:} Any qualified independent expert must be able to reproduce the exact same analytical findings using the same tools and documented methodologies.
    \item \textbf{Examination:} The report must survive to questioning by opposing counsel, who may deploy tactical strategies to exploit weaknesses.
\end{itemize}


\subsection{Rule 1: Define the Scope}
A well-defined scope protects the expert and prevents unnecessary disputes in court. The scope must explicitly cover:
\begin{itemize}
    \item \textbf{What Was Analysed:} Explicitly state the physical devices, datasets, file systems, target time ranges, user accounts, or specific operating systems fall within the scope of the assignment. Specific details such as hardware model numbers, serial numbers, date ranges, and complete data volumes should all be recorded.
    \item \textbf{What Was NOT Analysed:} This may be because the items fell outside the terms of the legal mandate, were technically inaccessible or were completely irrelevant to the specific legal question posed.
    \item \textbf{Why:} Every scoping decision must be justified.
\end{itemize}


\subsection{Rule 2: Declare Constraints and Limitations}
A report tmust contains also limitations (hardware fails, datapartially overwritten, ecc...)
\medskip

\noindent
Limitations must be categorized into three specific areas:
\begin{itemize}
    \item \textbf{Technical Limitations:} When some data is unrecoverable
    \item \textbf{Limitations of the Mandate:} The expert's task is strictly defined by the technical questions posed by the court or the appointing party. Analysis beyond those specific questions is not only unnecessary, but it may also be procedurally improper. Clarify what the mandate authorized and what it did not permit.
    \item \textbf{Limitations of Available Data:} Not all relevant data may have been made available to the investigator.\\
    For instance, logs may have been overwritten by normal activity 
\end{itemize}

\subsection{Rule 3: Guarantee Replicability}
Given access to the same source data and following the exact same documented methodology, an external expert must be able to replicate the same results. 
\medskip

\noindent
To satisfy this requirement, the reporting workflow must document:
\begin{itemize}
    \item \textbf{Forensic Copies:} All analysis must be conducted on verified forensic images (bit-for-bit copies) of the original storage media.\\
    Original evidenceaa must remain untouched and safely preserved.

    \item \textbf{Replicable:} Include documenting the exact sequence of operations, tools, the precise parameters applied, and any specific configuration decisions made during the analysis.
\end{itemize}

\subsection{Rule 4: Manage the Chain of Custody}
The chain of custody lifecycle is managed across five continuous sequential stages:
\begin{enumerate}
    \item \textbf{Seizure:} Recording the exact date, time, physical location, personnel present, and the serial number/condition of the device at the moment of collection.
    \item \textbf{Acquisition:} The creation of a forensic bit-stream image accompanied by immediate cryptographic hash verification.
    \item \textbf{Storage \& Transfer:}  Every transfer must be logged with formal signatures, timestamps, and condition verifications.
    \item \textbf{Analysis:} Only work on verified copies, log all software actions, and confirm the working copy's hash matches the original before starting.
    \item \textbf{Court Presentation:} Presenting the physical and digital evidence in court alongside the complete, uncompromised custody record to defend against any challenges to admissibility.
\end{enumerate}

\subsection{Rule 5: Uniquely Identify All Evidence Items}
Every item must be recorded using the three pillars of evidence identification:
\begin{itemize}
    \item \textbf{Serial Number} 
    \item \textbf{Cryptographic Hash Values} 
    \item \textbf{Image} The specific forensic image format used (e.g., E01, AFF, RAW), including metadata embedded at creation.
\end{itemize}

\subsection{Rule 6: Declare Tools and Methodology}

An expert report must explicitly detail:
\begin{itemize}
    \item \textbf{Software Used} 
    \item \textbf{Version Numbers} 
    \item \textbf{Configuration and Parameters} 
\end{itemize}


\subsection{Rule 7: Separate Facts from Interpretations}

The presentation of information must follow a clear, separate path:
\begin{enumerate}
    \item \textbf{Raw Technical Data:}These are verifiable facts backed by direct source references.
    \item \textbf{Analysis:} Where the expert identifies context, connects file system patterns, resolves conflicting timestamps, and provides clear attribution.
    \item \textbf{Conclusion:} 
    The conclusions must be expressed using clear, appropriately probabilistic language and must never stand as unverified assumptions.
\end{enumerate}
For example, writing \textit{"the user accessed the file"} is an interpretation; the raw fact is that \textit{"a file access event was recorded in the MFT at a specific timestamp"}
\subsection{Rule 8: Construct a Coherent Timeline}

Building a valid forensic timeline involves managing three distinct core challenges:
\begin{itemize}
    \item \textbf{Technical Events:} Gathering raw timestamps from multiple file system properties (created, modified, accessed, changed dates), Windows Registry activity, browser histories, metadata attributes, and system events (such as Syslog or Windows Event Logs).
    \item \textbf{Logical Sequence and Narrative:} Once events are mapped, the expert must establish a clear logical flow.
    \item \textbf{Handling Gaps and Conflicts:} Real-world timelines frequently contain missing records, altered fields, or conflicting timestamps. These gaps must be explicitly addressed rather than hidden
\end{itemize}
Investigators must clearly document the reference clock zone.


\subsection{Rule 9: Use Probabilistic Language}
Use clear probabilistic wording to express confidence levels and avoid absolute statements.

\subsection{Rule 10: Write for the Judge}
To maintain absolute clarity without compromising technical precision, the document must be compliant with three strict presentation standards:
\begin{itemize}
    \item \textbf{Clarity Above All:} Every technical concept introduced in the report must be clearly defined the first time it appears..
    \item \textbf{Logical Structure:} The report must follow a clear logical architecture.
    \item \textbf{Accessible Language:} Write using short direct sentences.
\end{itemize}
